% 第8章 功率放大电路
\chapter{功率放大电路}

前几章的放大电路都以“不失真地放大电压”为目标，输出功率很小；本章的\textbf{功率放大电路}（简称功放）要同时送出大电压和大电流，即在负载上得到足够大的\textbf{功率}。这一目标决定了它与电压放大电路在三方面根本不同：管子在大信号下运用（要用图解法、要考虑极限运用）、效率成为核心指标（不能把电源功率白白变成管耗）、非线性失真必须专门处理。

本章的结构是：先讲功放的特点与四种工作状态，再讲应用最广的乙类互补对称电路（OCL/OTL）的工作原理与功率、效率计算，然后是交越失真的消除、准互补与集成功放，最后是功放管的选择。互补输出级的具体电路形式在第3章（直接耦合互补输出级）已出现过，本章把它当作“功放”重新定量分析，不再重复其集成运放内部的偏置细节。

\section{功率放大电路的特点与分类}
\subsection{与电压放大电路的区别}
\begin{itemize}
    \item \textbf{追求的目标不同}：电压放大电路要求电压增益大、失真小；功放要求输出功率大、效率高。
    \item \textbf{信号幅度不同}：功放的输入、输出都是大信号，管子在整个动态范围内摆动，可能进入饱和区或截止区，因此\textbf{不能再用小信号 $h$ 参数等效电路}，只能用图解法（交流负载线）分析。
    \item \textbf{分析方法不同}：电压放大看增益、输入输出电阻、频响；功放看输出功率、效率、非线性失真度、最大输出功率，还要校核管子的极限参数与散热。
\end{itemize}

\subsection{主要性能指标}
\begin{equation}
P_{om}=\frac{U_{om}}{\sqrt2}\cdot\frac{I_{om}}{\sqrt2}=\frac{U_{om}I_{om}}{2}
=\frac{U_{om}^2}{2R_L}
\end{equation}
其中 $U_{om}$、$I_{om}$ 为输出电压、电流的\textbf{幅值}（不是有效值）。效率定义为输出功率与电源提供的直流功率之比：
\begin{equation}
\eta=\frac{P_{om}}{P_V}\times100\%,\qquad P_V=\frac{1}{T}\int_0^{T}U_{CC}\,i_C\,\mathrm{d}t
\end{equation}
此外还有非线性失真度 $\gamma$（输出中谐波总量与基波之比）、最大输出功率、频率响应、信噪比等。功放的三个指标互相牵制：\textbf{输出功率越大、效率越高，往往失真越大}，设计就是在这三者间取折中。

\subsection{工作状态的分类}
按管子在一个信号周期内的导通时间（导通角 $\theta$）划分，这是本章最重要的一张分类表：
\begin{itemize}
    \item \textbf{甲类（A 类）}：$\theta=360^\circ$，管子全程导通。静态工作点设在负载线中部，波形失真最小，但静态电流大、效率最低（变压器耦合时理论极限 $50\%$，电阻负载时更低）；单管甲类功放只用于小功率或要求低失真的场合。
    \item \textbf{乙类（B 类）}：$\theta=180^\circ$，管子只在半个周期导通。静态电流为零，效率最高（理论极限 $78.5\%$），但输出波形在正负半周交接处产生\textbf{交越失真}，且单管不能完整复现波形，必须用两只管子互补工作。
    \item \textbf{甲乙类（AB 类）}：$180^\circ<\theta<360^\circ$，静态时给管子一点正向偏置使其微导通。它兼顾了乙类的高效率与甲类的低失真，是实际功放最常用的状态，效率约 $60\%\sim70\%$。
    \item \textbf{丙类（C 类）}：$\theta<180^\circ$，管子导通时间不到半个周期，效率可高于乙类，但失真极大，只能用于负载为谐振回路的\textbf{高频调谐功率放大器}（谐振回路把窄脉冲中的基波滤出来），不能用于音频放大。
\end{itemize}

\begin{table}[htbp]
\centering
\caption{四种工作状态的比较}
\begin{tabular}{@{}lllll@{}}
\hline
状态 & 导通角 $\theta$ & 静态电流 & 效率（理想） & 主要问题 \\
\hline
甲类 & $360^\circ$ & 大 & $\le50\%$ & 效率低、管耗大 \\
乙类 & $180^\circ$ & 零 & $78.5\%$ & 交越失真 \\
甲乙类 & $180^\circ\sim360^\circ$ & 小（微导通） & $60\%\sim70\%$ & 需精确设置偏置 \\
丙类 & $<180^\circ$ & 零 & 更高 & 失真大，只用于谐振负载 \\
\hline
\end{tabular}
\end{table}

\formulaNote{易错}：$78.5\%$ 是\textbf{乙类互补}电路在输出电压幅度达到电源电压时的理论极限，实际因饱和压降、偏置电流、驱动损耗总要低一些（常见 $60\%$ 左右）。不要把它当成“乙类一定能达到”的指标。

\section{互补对称功率放大电路}
\subsection{乙类互补对称电路（OCL）的组成与工作原理}
电路采用\textbf{双电源} $\pm V_{CC}$，两只特性对称的管子 $T_1$（NPN）、$T_2$（PNP）接成射极输出器形式，负载 $R_L$ 接在两管发射极与地之间。输入信号同时加到两管基极，输出从发射极取出，因此电压增益近似为 1，靠电流放大提供功率。

工作过程：$u_i>0$ 时 $T_1$ 导通、$T_2$ 截止，电流由 $+V_{CC}$ 经 $T_1$ 流过 $R_L$；$u_i<0$ 时 $T_1$ 截止、$T_2$ 导通，电流由 $R_L$ 经 $T_2$ 流回 $-V_{CC}$。两管\textbf{轮流导通半个周期}，在负载上合成完整的正弦波，故称互补对称电路。由于没有输出电容、没有输出变压器，故称 OCL（Output Capacitor-Less）。

两个要点：
\begin{itemize}
    \item 两管都是射极输出器，属电压跟随，因此 $u_o\approx u_i$、输出电阻小、带负载能力强；
    \item 静态时两管都截止，输出端电位为零，负载上无直流电流，不需要隔直电容——这是双电源换来的好处。
\end{itemize}

\subsection{输出功率、效率与管耗}
设输出电压幅值为 $U_{om}$，则输出电流幅值 $I_{om}=U_{om}/R_L$。忽略管子的饱和压降，最大输出幅度可达电源电压，即 $U_{om}\approx V_{CC}$。

\subsubsection{输出功率}
\begin{equation}
P_{om}=\frac{U_{om}^2}{2R_L},\qquad
P_{om(\max)}=\frac{V_{CC}^2}{2R_L}
\end{equation}

\subsubsection{电源提供的功率}
每个电源只在半个周期内供电，一个周期内电源送出的总电荷对应平均电流 $2I_{om}/\pi$，故两电源提供的总平均功率为
\begin{equation}
P_V=\frac{2}{\pi}\cdot\frac{V_{CC}U_{om}}{R_L}
\end{equation}

\subsubsection{效率}
\begin{equation}
\eta=\frac{P_{om}}{P_V}
=\frac{U_{om}^2/(2R_L)}{2V_{CC}U_{om}/(\pi R_L)}
=\frac{\pi}{4}\cdot\frac{U_{om}}{V_{CC}}
\end{equation}
当 $U_{om}\to V_{CC}$ 时取最大值
\begin{equation}
\eta_{\max}=\frac{\pi}{4}=78.5\%
\end{equation}
值得注意：\textbf{输出幅度越小，效率越低}（$U_{om}$ 减小时 $\eta$ 线性下降），而管耗反而在中等输出时最大，这一点下面马上看到。

\subsubsection{管耗}
两管总管耗等于电源功率减去输出功率：
\begin{equation}
P_T=P_V-P_{om}=\frac{2V_{CC}U_{om}}{\pi R_L}-\frac{U_{om}^2}{2R_L}
\end{equation}
每管各分担一半。对 $U_{om}$ 求导并令其为零，得管耗最大时的输出幅度
\begin{equation}
U_{om}=\frac{2V_{CC}}{\pi}\approx0.636V_{CC}
\end{equation}
代回得每管最大管耗
\begin{equation}
P_{T\max}=\frac{V_{CC}^2}{\pi^2R_L}\approx0.2P_{om(\max)}
\end{equation}
这条结论非常重要：\textbf{最大管耗不出现在最大输出时，而出现在输出幅度约为最大值的 $64\%$ 处}，且每管最大管耗只有最大输出功率的 $1/5$ 左右。选管与设计散热器时必须按 $P_{T\max}$（而不是按 $P_{om}$）来校核。

\subsection{功率管的选择}
三组极限参数都要留有余量：
\begin{equation}
I_{CM}>I_{om(\max)}=\frac{V_{CC}}{R_L},\qquad
U_{(BR)CEO}>2V_{CC},\qquad
P_{CM}>0.2P_{om(\max)}
\end{equation}
第二条最容易被忽略：一管导通到饱和时，另一管的集电极—发射极之间承受的电压是 $2V_{CC}$（一端为 $+V_{CC}$、另一端为 $-V_{CC}$），所以耐压要按两倍电源电压选。此外还要考虑二次击穿与安全工作区（SOA），大功率管必须装散热器，按热阻与最高结温降额使用。

\subsection{单电源互补对称电路（OTL）}
OCL 需要双电源，便携设备无法提供。OTL（Output Transformer-Less）用\textbf{单电源}加一只大容量输出电容 $C$ 代替负电源：静态时输出端电位被 $C$ 充到 $V_{CC}/2$，电容相当于一个“$V_{CC}/2$ 的电池”。

工作过程：$u_i>0$ 时上管导通，电源经上管给负载供电并给 $C$ 充电；$u_i<0$ 时上管截止、下管导通，$C$ 通过下管向负载放电，充当负半周的电源。因此输出幅度最大只能到 $V_{CC}/2$：
\begin{equation}
U_{om(\max)}=\frac{V_{CC}}{2}-U_{CES}
\end{equation}
相应的功率与效率为
\begin{equation}
P_{om(\max)}=\frac{(V_{CC}/2)^2}{2R_L}=\frac{V_{CC}^2}{8R_L},\qquad
P_V=\frac{V_{CC}U_{om}}{\pi R_L},\qquad
\eta=\frac{\pi U_{om}}{2V_{CC}}
\end{equation}
$\eta$ 的极限仍是 $78.5\%$。输出电容 $C$ 的容量要足够大，使其在最低工作频率下的容抗远小于 $R_L$，否则低频段输出幅度下降、产生附加相移（这是 OTL 低频响应不如 OCL 的原因）。若把两路 OTL 反相驱动同一负载（BTL 桥式），输出幅度加倍，最大输出功率可提高到单路的 4 倍，这是低压供电下提高功率的常用办法。

\subsection{交越失真及其消除}
\subsubsection{交越失真的产生}
乙类互补电路的静态工作点为零，而管子有导通电压 $U_{on}$（硅管约 $0.6\sim0.7\ \mathrm{V}$）。当 $|u_i|<U_{on}$ 时，两只管子都截止，负载上没有电流，输出电压为零。于是正负半周交接处出现一段“死区”，波形在过零附近被削平，这种失真称\textbf{交越失真}。它与饱和失真、截止失真的区别在于：它发生在\textbf{过零点}，而不是波峰或波谷。

\subsubsection{消除方法}
思路是把电路从乙类改为\textbf{甲乙类}：给两管基极之间加一个直流电压 $U_{B1B2}$，使静态时两管都已微导通，交界处的死区被“填上”。工程上有三种做法：
\begin{enumerate}
    \item \textbf{二极管偏置}：在两管基极之间串入两只二极管（或用一只二极管加电阻），由电源经电阻供给直流电流，使
    \begin{equation}
    U_{B1B2}=U_{D1}+U_{D2}\approx1.2\sim1.4\ \mathrm{V}
    \end{equation}
    二极管动态电阻很小，交流信号几乎全部加在两管基极上（$u_{b1}=u_{b2}=u_i$），不损失增益；同时二极管的正向压降随温度下降，正好补偿管子 $U_{BE}$ 的负温度系数。
    \item \textbf{$U_{BE}$ 倍增电路}：用一只晶体管和两只电阻构成“恒压源”，其两端电压为
    \begin{equation}
    U_{B1B2}=U_{CE}=\frac{R_3+R_4}{R_4}U_{BE}
    =\left(1+\frac{R_3}{R_4}\right)U_{BE}
    \end{equation}
    调节 $R_3/R_4$ 可得到任意倍数的直流偏压，且温度系数可调，是集成功放中最常用的偏置方式。
    \item \textbf{恒压源偏置}：用稳压管或恒流源加电阻提供固定偏压，简单但温度补偿差。
\end{enumerate}
偏置电流要选得合适：太小则交越失真消除不净，太大则静态管耗上升、效率下降，一般取几毫安到几十毫安（视功率管大小）。

\section{准互补输出级与复合管}
\subsubsection{为什么要“准”互补}
理想互补要求 NPN 与 PNP 两管的特性完全对称。但大功率管中，PNP 型（尤其是硅 PNP）的制造工艺使其 $\beta$、频率特性、饱和压降都明显劣于同功率的 NPN 管，很难配对。解决方法是\textbf{都用 NPN 型大功率管作输出管}，靠复合管来“造出”一个 PNP 的等效管。

复合管（达林顿接法）把两只管子的电极按电流方向串接，等效为一只管子：
\begin{equation}
\beta\approx\beta_1\beta_2,\qquad
r_{be}\approx r_{be1}+(1+\beta_1)r_{be2}
\end{equation}
复合管的类型由\textbf{前一只管子}决定（NPN 驱动 + NPN 输出 = NPN 型；PNP 驱动 + NPN 输出 = PNP 型）。于是上管用“$NPN+NPN$”复合、下管用“$PNP+NPN$”复合，两只输出管同为 NPN 大功率管，特性易于配对——这种电路称\textbf{准互补输出级}。

准互补的优点是：输出管同型号便于配对、电流放大系数大（前级只需很小的驱动电流）、便于集成的偏置与保护。缺点是复合管的饱和压降较大（至少一个 $U_{BE}$ 加一个 $U_{CES}$），使最大输出幅度略有下降；且复合管容易因内部无泄放路径而关断变慢，实际电路中常在基极—发射极之间接泄放电阻。

\section{集成功率放大器}
集成功放把\textbf{前置放大级、推动级、互补输出级、偏置电路与保护电路}做在同一芯片上，外部只需接少量电阻电容与散热片。其内部结构与分立电路一一对应：
\begin{itemize}
    \item \textbf{输入级}：差分放大或射极跟随，提供高输入阻抗；
    \item \textbf{偏置级}：$U_{BE}$ 倍增电路，设定甲乙类静态工作点，消除交越失真；
    \item \textbf{输出级}：准互补或互补对称（多为 OTL 形式，单电源供电、外接输出电容）；
    \item \textbf{保护电路}：过流保护（检测输出电流、限制驱动）、过热保护（芯片温度超过限值时自动降低输出）、短路保护。
\end{itemize}
典型产品如 LM386 一类低压小功率音频功放：单电源 $4\sim12\ \mathrm{V}$ 供电、输出功率几百毫瓦、电压增益可由外接阻容在 $20\sim200$ 之间设定，用于收音机、小音箱。使用时必须注意三点：电源端就近接去耦电容（防自激）；输出端对地的 RC 网络（茹贝尔网络）抵消扬声器的感性负载，防止高频自激；芯片功耗接近额定值时加足够大的散热片——集成功放的损坏绝大多数是过热或输出短路造成的。

\section{功率放大电路的读图与估算步骤}
遇到一个功放电路，按下面顺序估算，不容易乱：
\begin{enumerate}
    \item 看电源：单电源（OTL，输出端有电容、静态电位 $V_{CC}/2$）还是双电源（OCL，输出端静态电位为零）；
    \item 定最大输出幅度：OCL 为 $V_{CC}-U_{CES}$，OTL 为 $V_{CC}/2-U_{CES}$；
    \item 算最大输出功率 $P_{om(\max)}=U_{om(\max)}^2/(2R_L)$；
    \item 若给出实际输出幅度，先算 $P_{om}=U_{om}^2/(2R_L)$，再代入 $P_V$ 求效率 $\eta=P_{om}/P_V$；
    \item 校核管子：$I_{CM}$、$U_{(BR)CEO}$、$P_{CM}$ 三项与前面的选管条件对照，并检查散热；
    \item 判断工作状态：静态有无偏置（有则甲乙类）、输出波形过零处是否削平（有则交越失真未消除）。
\end{enumerate}

\formulaNote{易错}：估算时“电压用幅值还是有效值”最容易错。$P=U^2/(2R_L)$ 里的 $U$ 是\textbf{幅值}；若用有效值则 $P=U^2/R_L$。$P_V$ 中的 $2/\pi$ 系数来自“半波导通取平均”，OTL 只有一个电源，系数是 $1/\pi$，不要照抄 OCL 的式子。

\section{本章小结与易错点}
\begin{itemize}
    \item 功放与电压放大的本质区别是“大信号 + 追求功率与效率”，因此分析工具从等效电路换成图解法，指标从增益换成 $P_{om}$、$\eta$、$\gamma$。
    \item 四种工作状态按导通角分：甲类 $360^\circ$（效率 $\le50\%$）、乙类 $180^\circ$（$78.5\%$）、甲乙类介于两者之间（实际 $60\%\sim70\%$）、丙类 $<180^\circ$（只用于谐振负载）。
    \item OCL 双电源、无输出电容、静态输出为零；OTL 单电源、输出电容充当负半周电源、静态输出为 $V_{CC}/2$。两者的最大效率都是 $78.5\%$，但最大输出功率公式差 4 倍（$V_{CC}^2/2R_L$ 对 $V_{CC}^2/8R_L$）。
    \item 最大管耗出现在 $U_{om}\approx0.636V_{CC}$ 处，每管 $P_{T\max}\approx0.2P_{om(\max)}$；选管耐压要按 $2V_{CC}$（OCL）。
    \item 交越失真发生在过零点，根源是乙类无静态偏置；消除办法是加二极管偏置、$U_{BE}$ 倍增电路或恒压源，把电路变成甲乙类。
    \item 准互补解决的是“大功率 PNP 管难配对”的工艺问题，关键是复合管的类型由前一只管子决定。
    \item 集成功放的内部结构与分立电路一一对应，外部关键是去耦、茹贝尔网络与散热。
\end{itemize}
